
% 来源：EPUB part063.xhtml
\section{好公司最重要}

就投资而言，好公司最重要。（2015-11-01）

% 来源：EPUB part064.xhtml
\subsection{想10年20年后}

当你迷惘的时候，试试往远处看？看10年往往比看1到2年要容易很多，看一两年会比看一两天容易很多。比如投资上，很多公司其实是很容易知道它们10年后很可能会不行的，那为什么现在要买呢？避开这些公司你会发现其实你的选择会特别少，但心里会踏实很多。（2022-01-26）

\BookQuestion{网友：A股除了茅台，还有哪个公司你认为生意模式不错？}

作为一个懒人，我能了解的公司很少。不过你这个问题翻译过来就是10年后（或者20年后）的某一天来看现在我应该用茅台来换哪个公司呢？假设一个礼拜后股市就关闭了，10年或20年后的某一天再开，你的决策会是什么？（2019-09-08）

大多数东西10年后啥都不是，但确实有些东西10年后会变化很大，还有些东西几十年后和几十年前一样。自己去找吧，人生的酸甜苦辣咸都在里面了。（2022-10-30）

\BookQuestion{网友：如果大道发现低估的股票，会不会去弄明白市场哪里错了？我们假设市场大多数时候是对的吧？如果不知道市场错在哪，是不是意味着我们还没真懂？}

想公司，想生意，想10年20年后公司是什么。不要理会市场，你永远不会明白疯子在想什么的。（2024-03-14）

贵或者便宜都是从未来10年或更远的角度看的。比如，如果你觉得一家公司未来肯定会完蛋，那么现在什么价格都是贵的。几年前有个某公司的第二大股东聊到该不该卖某公司，什么价钱卖的问题。我说，这个看着早晚为0的公司，什么价卖都是好价钱吧？尽量避开那些看起来10年后日子会不好过的公司很重要，偶尔错失一些机会不会伤害到你的。（2020-12-06）

\BookQuestion{网友：想想东方甄选自营品10年20年后的样子……是不是很有可能成为中国的线上山姆、开市客？}

想象力丰富一般不是坏事，但对投资而言则未必是好事。（2024-03-14）

长线持有一定是在买对好股票的前提下。有很多人往往只是在“被套”的时候才想起“长线”来。

这里“被套”的意思是指有些人本来只打算拿10天的股票，结果因为“亏钱”了而决定改为拿10年，这种决定往往会让人亏更多。不打算拿10年的股票一定别去打算拿10天，不然就是“炒”了。（2011-12-13）

你认为10年后还在的公司都有谁？定义是10年后营业额和利润都不低于今天的才可以认为还在哈。

我觉得苹果、茅台、腾讯应该都还在；谷歌、微软、伯克希尔应该也还在。我希望OPPO和vivo还在，但不是100\%确定。

20年呢？我觉得苹果、茅台应该还在，谷歌、微软、伯克希尔也应该还在，腾讯90\%还在，OPPO和vivo应该75\%还在。

30年呢？……确属脑力游戏，没事想想这个也许能对投资有点帮助。（2024-01-28）

在年利润10亿美元以上或年利润50亿人民币以上的公司里，有谁未来10年的总利润会比过去10年高？在这些企业中，有谁未来10-20年的获利总和会比未来10年的总和高？

\BookQuestion{网友：大道是在寻找投资标的吗？}

这是我的思维方式，我大概就是这么看公司的。当你这么看公司时，很自然就重视商业模式了。

我了解的公司很少，在我比较了解的公司里（年利润10亿美元以上的），我觉得未来10年总利润比过去10年总利润高的公司可能有苹果、茅台、腾讯、伯克希尔、微软、谷歌，拼多多应该算一个，但那是因为拼多多过去10年还没怎么赚到钱。麦当劳、可口可乐这类公司打败通胀的概率应该很大，所以也应该算。有点汗颜，我发现我了解的公司实在是太少了。（2024-01-29）

\BookQuestion{网友：四大国有银行10年后还会存在，所以银行是一门好生意吗？}

日本甚至有存在了上千年的生意，你可能完全没听说过，但你可以肯定那不是什么好生意。生存时间长和是不是好生意没有必然联系。（2024-03-14）

% 来源：EPUB part065.xhtml
\subsection{好公司不怕萧条}

好公司不怕萧条，大萧条来的时候好公司的额外优势更大，公司价值所代表的实际购买力反而会相对上升。所以，投资要聚焦在找好公司上面。（2019-09-19）

我喜欢孙正义说过的一句话：越是迷惘的时候，越是要往远处看。任何时候好公司都是值得买的，尤其是大环境那个啥的时候。（2024-10-16）

从投资的角度讲，找到好公司比啥都重要。在大家有“危机”的时候，好公司往往能将“危机”简化为“机”的。（2018-09-05）

\BookQuestion{网友：感觉实体企业的寒冬真的来了！}

谁的冬天？关于冬天的说法，茅台肯定是不同意的，我们也不太同意，有类似观点的应该还有很多家企业，比如苹果、谷歌……（2016-02-23）

\BookQuestion{网友：您说“这些年花在投资上的时间不多，但投资水平有很大长进”，具体是什么意思？}

就是不再受市场影响，更加坚信要买那些有好的生意模式和好的企业文化的好公司。（2019-07-28）

比如茅台，过去30年里不管你在酒前酒后、左侧右侧买，10年后看都是好投资，但用margin买的还是有可能要裸奔的。所以，不懂不碰，不用margin就是不为清单。（2020-10-11）

\BookQuestion{网友：巴菲特说在认识芒格之前，他特别喜欢在市场上捡便宜货，是芒格让他意识到花贵一点的价格买个好货回报高过捡便宜货。我觉得巴菲特那时候之所以喜欢便宜货应该是因为对商业的深层次理解等都还不够，对贵货很难下手。价值投资需要对行业、公司、商业，甚至对社会有很深的理解。我不认为年轻人容易具备这种能力。相反，便宜货会比较容易判断些。}

呵呵，还是个未来现金流折现的问题。巴菲特发现“便宜货”的未来现金流折现比有把握成长的公司的未来现金流折现来得还要小，这可能就是芒格提醒的作用吧。但是，看懂好的成长公司可是要比捡便宜难得多哦。（2010-05-24）

\BookQuestion{网友：总觉得自己在平庸公司里徘徊不前，如何才能跳出怪圈？}

关键是要搞懂生意，能看长远。“时间是平庸公司的敌人，是伟大公司的朋友。”这句话也许能帮到你。（2020-12-02）

找到好公司确实很难。所以，真要投资的话，要坚持的是要去找好公司，去看懂生意。至少要学会坚持不买烂生意。（2024-08-15）

“买自己一眼就看懂的生意往往都不错”，说这句话的人很有悟性哦。巴菲特也讲过，伟大的生意一般只要一个段落就可以说得很清楚。（2013-02-01）

其实愿意化繁为简的人很少，大家往往觉得那样显得没水平，就像买茅台一样，没啥意思。（2019-08-14）

复杂的事情简单做，简单的事情重复做。（2023-05-29）

\BookQuestion{网友：20世纪我们看到了无数伟大的公司，包括沃尔玛、IBM、微软。到21世纪，进入了信息时代，什么样的公司会最受欢迎？}

很多20世纪的伟大公司今天依然是伟大公司。有很多公司的实力也是难以想象的。刚刚不小心看到思科（CSCO）居然有390亿美元现金，可以单独去救欧洲了。（2010-05-15）

“如果你因为某样东西被低估而买入，那么当它价格接近你所计算的内在价值时，你就得考虑卖出。这很难。但如果你买入几家伟大的公司，那你就可以坐享其成。这可是件美事。”（芒格）

这条大概要花很多年才能真明白，大概就是未来现金流折现的意思。我觉得自己现在大概有点明白了。（2012-06-26）

\BookQuestion{网友：有人说资金量小可以关注成长性更强的小公司，您的建议呢？}

我没啥建议，你要不懂投资，其实都是一样的，你要懂投资其实也没啥区别。（2022-01-23）

投资的原理是一样的。找小公司也是要找那些商业模式好、企业文化好，但规模和市值都比较小，大家还没发现的那种。现在回过头来想，多数小额投资者投小公司有可能情况会更糟，因为小公司往往资料比较少，小额投资者更容易中招。但假如我自己作为一个小额投资者（依然比较懂企业），如果有时间的话，我会努力去找那些我认为大有前景的小公司，这样回报会高不少。（2019-04-30）

\BookQuestion{网友：资金不大时是捡烟蒂还是选择优秀的公司？}

对绝大多数人而言，其实和资金大小一点关系都没有。（2013-12-18）

小资金的人对企业的理解一般也容易弱些，不容易看懂伟大、优秀企业与普通企业的区别，所以多数人都是容易从烟蒂开始的。只有很少的人最后会真正悟到投资的道理，并同时具备看懂伟大企业的能力。（2015-01-19）

“小资金的人如果懂投资的话是可以找到更好的机会的”，这其实是个伪命题，因为小资金能懂投资的概率也是小的。（2025-03-05）

\BookQuestion{网友：购买市值低于净资产，甚至是净现金的公司，特别要注意哪些问题？}

我觉得你需要对公司特别了解才行。这种情况特别容易陷于所谓的“价值陷阱”，就是看着便宜但永远也不涨的那种。如果从未来现金流折现的角度看则会容易理解些。（2010-10-21）

市值低于现金的公司千万要小心（和投资别的公司一样小心就好）！一眼看上去市值低于现金的公司确实经常是值得花时间去看的，但多数情况下这未必就是便宜货。（2018-01-12）

\BookQuestion{网友：可以确认的净资产是否就是股价向下空间的一条底线，一旦低于这个底线，就应该大量买入？}

我的回答还是不一定。比如100块的东西，如果98块卖给你，也不一定合算，除非你马上倒手一下。（2010-04-27）

\BookQuestion{网友：面对成千上万非上市的小公司，除了观其行听其言、深入了解其产品和服务，还有什么办法能高效率找到未来有潜力的小公司？}

不知道。我只关心上市公司，因为容易了解很多。没有10年8年的可追溯历史，很难了解一个公司，了解一个小公司几乎是个不可完成的任务，除非开公司的人刚好是你以前就认识了很久的人。（2023-04-26）

% 来源：EPUB part066.xhtml
\subsection{越是懂投资越应该集中}

那些买了一点点的不算看懂了，如果真的看懂了却只买了一点点是说不过去的。

我对看懂的定义非常简单，就是敢下重注。按巴菲特最多分散到6个公司的说法（我很同意哈），看懂并敢下注的比例至少要占1/6吧？不然都应该叫“观察仓”？

这里不是说一定要达到1/6的意思，因为有时候价格（机会）并不是那么合适。（2024-04-21）

\BookQuestion{网友：您觉得个人投资者是集中投资好，还是适度分散好？}

跟是不是个人投资者无关。我觉得越是懂投资的越应该集中，但据说不太懂投资的人则最好分散到指数类基金，不过我没试过。（2020-10-10）

\BookQuestion{网友：一开始做投资没有想得那么明白很正常，似懂非懂的时候拿点小钱试一试也未尝不可。}

多数人的问题是算得太多而想得太少。拿点小钱试试没啥不可，但你想试的是什么？市场的走势？（2018-01-12）

\BookQuestion{网友：茅台和苹果的资金分配可以量化么？比如这个3成那个5成？}

从来没算过，也不觉得重要。（2019-05-27）

我一生中单吊一股不是一次两次了，创立公司是，投网易也是。巴菲特也不是一次两次了。（2023-05-12）

% 来源：EPUB part067.xhtml
\subsection{每一次买卖都是独立的}

平常心看，每一次买卖其实都是独立的。（2013-07-04）

一次有个朋友飞来旧金山，我去机场接。由于早到，我在停车场投币器里投了一个小时的钱。结果他也早到了，所以我们回到车上时，投币器里还剩半个多小时的时间没用完。于是我们决定为了不浪费已经投进投币器里的钱，在停车场里的车上又待了半个多小时。

现实当中也许没人会真的在停车场等，但据说有个调查显示，500强CEO里有85\%的人承认他们干过类似的事情。本人属于85\%这类人。（2010-03-30）

呵呵，这个故事讲的是沉没成本。沉没成本是指由于过去的决策已经发生了，而不能由现在或将来的任何决策改变的成本。故事讲的就是为了去救已经沉没的钱又花了新的投资，结果是白花了。你可以想想自己有没有干过类似的事就明白了。（2010-04-02）

\BookQuestion{网友：引申到投资上，多数人被套住的时候都会守在那里，直到解套，而不是比较解套和其他标的哪一个机会更大。其实每天的收盘价就是我们的机会成本，没有卖就等于买了，否则可以买其他。所以非常认同段老师的卖和买入成本无关的说法，可惜能过这一关的人太少了，我自己也是挣扎了几年才爬过这一关的。}

赞。（2010-03-31）

讲个真实的小故事吧，大概差不多快10年前了，公司的同事或者说朋友们告诉我，弟兄们买了很多忠旺的股票，麻烦让我看下这家公司。于是我确实飞过去看了这家公司。回来后我告诉弟兄们：我觉得这家公司商业模式一般，企业文化一般，大家卖了吧，无论亏赚，换茅台或者苹果或者腾讯都应该更好啊。当时忠旺股价在3块左右。大部分人很快就换了，但确实有人不舍得换，总是说会有机会的会有机会的。一转眼这么多年过去了……机会成本杠杠的。（2020-04-26）

% 来源：EPUB part068.xhtml
\subsection{合适价钱就好}

其实我对好价钱已经不太提了，合适价钱就好。（2019-05-20）

好价和长期持有的关系是什么呢？其实是好股长期持有就是好价了，而不是买到了好价钱的意思。（2019-10-14）

\BookQuestion{网友：最近看了《穷查理宝典》，提到最佳的投资方式是以公平的价格买入伟大公司并持有。伟大的公司是等待市场定错价时再买入呢，还是见到公平价格就可买入？}

你一旦真的明白了什么是未来现金流折现，就可以解答你的问题。（2015-06-07）

\BookQuestion{网友：合适的价格有什么模糊的标准？}

价格合适主要指的是你自己的机会成本（也可以扩展到社会的机会成本）。比如你喜欢一个公司，你觉得未来20年的回报能有8\% /年，但同时你又是一个很会赚钱的人，平时总能赚到10\% /年，那你当然就可以考虑等等，看看有没有更好的机会。如果你只是把钱放在银行吃1\% /年的利息的话，年回报8\%就很好了，千万别听隔壁小马赚了很多钱的那些故事，哪怕那故事是真的。举个极端的例子，有个公司你觉得回报能有10\% /年，但银行利息是15\%，你该怎么办？（实际情况往往是股价会掉得很厉害。）（2019-05-17）

\BookQuestion{网友：你说过好公司价格不太重要，过得去就行，时间到了什么都会有的。请问背后的逻辑是什么？能以苹果或茅台为例子讲讲吗？}

用10年看就明白了。今天你看茅台，你觉得120元买的和180元买的有多大区别？当然，你如果知道茅台会到120元则不一样。很久以前我说过，苹果400美元很便宜，500美元也很便宜，到600美元时我依然说还是很便宜（在1/7之前），现在看是不是就容易明白了？（2019-09-12）

\BookQuestion{网友：茅台还没买够，但是已经涨起来一些了，安全边际降低，但是股息率依旧高于银行利息，这时候该如何操作？}

茅台以前还更便宜呢。尽量往远处看吧，别老想着赚点快钱。

\BookQuestion{网友：依稀记得大道2014年时说200元买入的茅台和100多元的茅台从长远看没什么区别，这话我都回味了近十年啦。}

很多人说100元当然比200元好，傻子才不知道呢！我们不知道的是当茅台180元的时候，还会掉到100元吗？我们知道的是180元的茅台已经很便宜，有钱就先买了。万一真掉到100元了，有闲钱再买，没有就跟自己无关了。（2024-09-27）

\BookQuestion{网友：既要选好股，又得择好时。}

号称自己会择时的人大多最后日子都不太好过。（2022-03-09）

“不择时”我的理解大概是：当我觉得从10年的角度看不贵而手里正好有钱的时候就可以买了，不需要等更便宜的价钱，因为这个价钱从10年的角度看已经很好了。（2023-02-23）

当合适的价钱出现时，除非你知道等什么，不然你在等什么？其实从长期的角度看，一点点价钱的差异远不如公司好坏来得重要。（2014-07-26）

我看到有些人手里拿着钱，看见别人赚钱心里着急。眼看着股价越涨越高，心里跟着越来越急，最后在“不择时”的鼓励下，一咬牙一跺脚就冲进去了。不择时不是追高的理由。真正不择时的人是真的看10年的人，而且人家钱大多早就已经在投资上了。追高不择时的人们实际上对短期的下跌是非常难以忍受的。（2023-02-23）

\BookQuestion{网友：如果看好的股票到达了可以接受的价格，您是习惯一次性全仓买入，还是分批买入？我很纠结。}

你纠结的东西不是关于价值的，10年后回头看没啥意义。除非你钱太多，不然见到好价钱只买一点点是件说不过去的事情。（2019-07-22）

我的“满仓主义”就是一般情况下，要投资的钱都放在投资上，手里不专门留准备用作投资的现金，因为现金长期来说是会输给通胀的。我的满仓主义最大的特点是不太在意“择时”，但非常在意“商业模式 + 企业文化”，当然前提是10年后看回来的价格是便宜的。满仓主义不意味着就一定要满仓，没有合适投资标的时（包括价格过高时）是可以拿着现金的。价格过高的意思是：这个公司的未来现金流折现达不到目前市值的水平。（这句话的潜台词是，我自己的认知范围里觉得这个公司赚不到那么多钱，所以价格就是过高的，不等于实际上这家公司做不到。）（2023-04-07）

% 来源：EPUB part069.xhtml
\subsection{持有 = 买入}

\BookQuestion{网友：如何理解持有 = 买入？我也是重仓苹果和茅台，但总是会害怕国际宏观环境。身边的人也总是提醒要控制仓位，后面跌了可以买更多。很希望能够有您这种持股不动的心态！}

理论上（忽略所有交易成本的前提下），持有 = 买入，或者买入 = 持有，其实不是一种观点，而是一个事实。我以前也为这个观点困惑过很长时间，直到有次跟巴菲特聊天时我问他，持有 = 买入吗？他说当然啊，这是一个事实啊！那之后我就突然觉得自己明白了。不过，明不明白这点，和投资好坏似乎也没太大关系。对你的苹果或茅台而言，不管你想买进、卖出或不动，和你已经持有多少、成本多少实际上都是无关的，是个独立的投资决策。（2020-10-12）

\BookQuestion{网友：持有等于买入，是否在讲“为卖出而买入（buy to sell）”和“为持有而买入（buy to keep）”的区别呢？}

这个东西没有标准答案，但我认为用buy to keep和buy to sell来区别是个非常好的说法。事实上，“持有”在每个时点确实是等于“买入”的，机会成本就是这个意思！然而，如果一个投资者真的那么想的话，就很容易掉到每时每刻都在想“是不是高估了”“是不是高估了”“是不是高估了”，然后过多地注意市场而不是聚焦在生意本身上。投资的本质在于生意本身，对市场的关注越高，投机的成分越高。所谓的价值投资者是应该可以完全不在乎某公司现在及未来是不是上市公司的因素而决定买卖该公司的。（2014-08-30）

\BookQuestion{网友：如果你持有三只股票，其中一只泡沫比较大，其他两只估值相对合理，你会调整仓位把泡沫的股票换到另外两只吗？}

我觉得难得的好公司最好还是拿着，换来换去比较复杂。前段时间在我管的一个朋友账号上把苹果换到伯克希尔了，换亏了不是（少赚了）？最近股市这个涨法我也看不懂，可能真的和利息太低了有关吧。我一直不太习惯买很高P/E的公司，但我们现在主要持有的苹果、茅台和腾讯P/E都是30多倍了（具体多少还没看），持有 = 买入啊，对我也是个考验哈。准备复习一下大道以前讲的东西，努力往前看10年20年。（2020-08-25）

这里表达的东西其实和苹果无关。昨天下午和一个朋友打球时，他告诉我他上午“又”买苹果了，大概560美元/股。记得上一次他买苹果是当苹果360美元的时候，结果买了以后就设了个399美元的卖单。我劝他别卖，因为这价太便宜，可他最后还是卖了。今天，其实苹果还是那个苹果，但市场价到过600多以后，似乎在500多下买单就容易多了，呵呵。

其实，卖股票和成本无关，买股票和其到过多少价也无关。最重要其实也是唯一重要的是公司本身。

什么时候确实可以完全忽略市场的影响的时候，对价值投资的理解就算是很到位了。

当然理解价值投资不等于就能在股市上赚到钱，就像知道要“做对的事情”的人不一定具备“把事情做对”的能力一样。

把事情做对需要有很多年的辛苦积累，不是看一两本书或者上几个“高手”的博客就能学会的。

大概由于这些天苹果掉得很厉害，昨天上午有个朋友打电话问我苹果发生什么事了。我说我不知道，但我重复了《功夫熊猫》里Master XiFu（就是广东话里师傅的意思，好莱坞的人大概以为中国人讲的都是广东话）的那句话：You must believe！

很多人读过很多书，知道很多投资的教条，But they just don't believe. Maybe they don't believe because they don't really understand. Who knows？（但他们就是不相信，可能他们不相信是因为他们并不真正懂得，谁知道呢？）

这里的“believe”千万不要理解成简单的自信，这个“believe”的建立是需要很多年的经验和理解的，但不是有很多年就一定可以理解的意思。我见过很多有很多年“经验”的人，其实从来就没有真正花时间去“理解”过。（2012-04-26）

\BookQuestion{网友：不买入 = 不持有吗？我在迷茫假如我持有一只股票，但是现在我又不想加仓了，是否干脆减仓甚至清仓？}

你的问题其实有点复杂的。持有确实等于买入，但未必就等于要加仓，因为往往你还要想看看是不是还有别的机会？不过，我自己经常就是有钱的时候就加仓了。网易从不到一块美元一直加到几十块美元（后来1分4了），苹果从300多美元（1分7之前）一直加到170-180美元，茅台前段时间500多元的那次还加了不少，都是因为刚好有现金流了。

其实从10年后看回来，这个问题没那么复杂。你迷惘的原因其实是因为害怕股价短期的变化。（2019-09-04）

\BookQuestion{网友：持有伟大公司的股权，在高估的时候，是该卖出还是继续持有？}

这是一个很难的问题，没有总是对的答案。巴菲特的观点是一点点高估不应该卖，几乎每次他卖了后都被证明是错的。我现在也慢慢理解巴菲特的观点了。找到一个伟大公司非常难，不要为一点点差价就卖掉了。（2019-03-20）

巴菲特在回答什么时候卖出伟大公司的时候说“never”！这个never不容易懂，巴菲特自己花了几十年才大致明白了。（2017-08-04）

\BookQuestion{网友：不能教条主义，散户还是遵循低估买入，高估卖出比较合适。}

非常对，只有这样才能维持散户的地位。（2019-06-22）

伟大的公司要惜卖！伟大的公司卖掉后大概率很难买回来，而且因为手里拿着现金，很可能会投到没那么好的生意，损失会有可能是双重的。（2022-01-31）

\BookQuestion{网友：大道面对不确定性的时候会焦虑不安吗？总觉得自己的决策和行动都达不到100\%，而且永远没可能达到100\%。在自己觉得能看懂的情况下，选择重仓后，会害怕吗？}

不用margin买了自己喜欢的生意，为什么会焦虑？苹果掉的时候我会焦虑吗？茅台掉了这么多我会焦虑吗？当你觉得焦虑的时候就说明你没搞懂你买的生意，那你就该离场了。买了不懂的生意是很容易焦虑的，也很难拿住。（2024-08-08）

真懂了就拿得住，不需要什么技巧。

再说一下关于拿住的观点：买股票（公司）和其曾经到过什么价无关，卖股票和买入价无关。不要老以为谁是高手，也不要老想着成为高手，老老实实找些自己能看懂的好公司，好价钱时买下拿着就好了。巴菲特讲过，最好的公司就是永远都不想卖的公司。一生找到几个这种公司，想不发达也难啊。（2015-01-13）

\BookQuestion{网友：您买卖股票的标准是什么？}

因为买公司是基于好的商业模式或企业文化，所以当觉得这两个方面长期来看出现不可修复的问题时，就要考虑离场了。（2022-03-03）

你觉得跟踪每天的新闻会对你有帮助吗？如果没有的话，那你为什么要这么做？如果你投资了一家企业后还会为每天的新闻改变你对企业的看法，那大概就说明你对这家企业其实不了解。当然，交易者（traders）是需要根据每天的新闻去做事情的，基本上，那样做应该会很累而且不太赚钱。另外，关心和跟踪是有点不太一样的。关心是一个享受的过程。

(2014-02-05)

% 来源：EPUB part070.xhtml
\subsection{根基不变没必要太关注宏观}

“我和沃伦并非依靠预测宏观经济走势并基于这些预测结果下注，才在人生中取得如今的成就。”（芒格）

我也不是。（2012-06-26）

“我对总体经济一窍不通，汇率与利率根本无法预测，好在我做分析与选择投资标的时根本不去理会它。”（巴菲特）

好多人以为自己看看新闻就能比巴菲特明白，可笑不？

“从预言中你可以得知许多预言者的信息，但对未来却所获无几。”（巴菲特）

哈哈，这个说法很妙。（2012-07-28）

\BookQuestion{网友：您研究宏观经济吗？还是把全部的精力放在研究微观经济上？}

对宏观经济完全没概念也是很难了解微观经济的。所谓对宏观经济的了解应该是对经济现象的理解，而不是花很多功夫去根据宏观经济的短期数据来不断调整自己的“投资”。我本身本来就学过些宏观经济，主要是马克思以及其他西方经济学的东西，所以对宏观经济知道一点。（2011-01-05）

其他的东西对公司短期盈利有影响，但如果你从5年10年的角度看，你会发现宏观的东西对公司的影响要比想象的小得多。（2010-02-07）

专注在公司生意上，忽略短期的波动很重要。眼前这些事10年后都不是事。（2018-09-10）

\BookQuestion{“2月28日，特朗普与泽连斯基在白宫发生争吵。”（新闻）}

终于有点明白巴菲特为什么突然从满仓主义者变得如此保守了。

\BookQuestion{网友：说巴菲特不关心宏观其实也是不太对的。}

在根基不变的情况下确实没必要关注太多的宏观变化。（2025-03-02）

% 来源：EPUB part071.xhtml
\subsection{为自己想拥有的股票投保}

卖put等于对投资回报的要求不像大家想的那么高，大概就是宁愿放弃一口吃个大胖子的机会，但更有机会有还过得去的回报，比如10\%/ 年，也不完全排除更高的收益。不过，卖put这个说法很危险。不是谁的put都可以卖的。我主要只是卖苹果的put。（2022-11-12）

我卖put的情况多是自己喜欢的公司的价钱在不高不低的区间，手里又正好有点闲钱。想着10年后这个价钱也不贵，闲着也是闲着，于是就卖个put，赚个年化十来个百分比的回报。特别便宜的时候最好别卖put。因为很可能赚了一点蝇头小利而失去大机会的。（2023-03-18）

我做股票的第一个交易就是卖put。一直不是很明白为什么大多数人搞不懂卖put和卖coveredcall（备兑看涨期权）的道理。照说如果明白买股票就是买公司的道理后，应该很容易明白卖put无非就是让你更便宜地买你原来想买的股票的办法而已。

\BookQuestion{网友：你已经很想买正股了，当然不会介意以更便宜的价格买。}

对，卖put的根本道理就这一条，不知道为什么会被人说得那么复杂。

\BookQuestion{网友：卖看跌期权本质上就是卖保险。当你对投资对象很有信心的时候，这真是一个好工具。}

我也一直用卖保险来比喻。但两者最本质的差别其实是保险的费率大概是可以算出来的，put不能。卖put实际上就是为自己不拥有的东西投保。我只为自己想拥有的股票投保。

\BookQuestion{网友：option（期权）可以作为价值投资的辅助工具？}

是啊，巴菲特也用啊。不过，巴菲特规模太大，只能卖点标普的put。我自己也觉得越来越难用了，规模大了以后市场没容量，卖不出好价钱就没意思了。（2011-12-18）

使用option时最重要的就是不要为option而option，只有对基本面非常明白时才有可能有时用一下。还有就是不要用margin做option。（2011-10-04）

我说的卖call一定指的是coveredcall，和市场无关但和了解有很大的关系，不了解的公司绝对不能碰任何option。卖put也不需要“大量”的资金，钱多多做些，没钱就不做，和别的一样。（2011-07-16）

从本质上来讲，卖option就是在赚投机者的钱，实际上是在和投机者对赌。投机者像来赌场的赌客，卖coveredcall和有些股票的put则有点像赌场。（2011-07-17）

你如果坚持用买call的办法去买你想买的公司（其实你只是想赌个大涨而已），最后的结果大概率是亏光光的。我自己是不买option的，不管是call还是put。（2020-11-13）

\BookQuestion{网友：如果有一家真正懂的公司，option是能增厚不少利润的。所以确定性最重要。}

能不能增厚利润不是很确定，但下手难度减少了不少，内心没那么折腾，受市场影响会降低一些。（2023-09-29）

卖put看着好玩，但背后的逻辑还是严密的。首先，从十年的角度看，170美元的苹果其实不贵（功夫全在这里），然后就好玩了。（2022-01-08）

\BookQuestion{网友：巴菲特谈期权时说如果你想要买卖股票，就应该直接买卖股票……}

巴菲特是对的！但这不意味着我的办法对我不对。卖put可能会少赚钱，但至少买入成本是低一些的。反正我每天卖1000个腾讯put会继续，比直接买股票要容易下手很多。顺便说下，我的大部分苹果就是通过卖put买回来的，虽然没巴菲特买得多，但比重可比巴菲特要高哈。

其实我的想法很简单，就是我管的钱一定要战胜通胀，不然我就亏钱了。所以，首先要找自己看得懂的好公司。好公司一般都有点贵，直接买老是觉得有点下不了手，但卖个put，就算最后买不到股票，年化十几个百分点的收入也还是可以的嘛。这里并没有一个所谓的卖option的策略，最重要的还是也只能是找到好公司。我并不追求所谓的最好的投资策略，我也不可能每天去办公室上班（我确实有个办公室，到目前为止只去看过一次）。不客气滴说，如果我全天候做投资，我的回报肯定可以高一些，但我为什么要这么做？我目前这个办法其实还是蛮有效的，一个人可以管人家两三百号人管的钱，而且好像过去20年表现都还不错哈（但并不总是能赢标普500）。绝大部分时间里我都是满仓或者实质性满仓的，最近有点闲钱，主要是有段时间的闲钱买了苹果同时卖了call，前段时间被call走了一部分。目前价位下，买苹果真有点下不去手，但卖腾讯的put似乎心里还是踏实的。苹果以前回调过很多次啊。也许最近大盘涨得有点邪乎，感觉怎么懂苹果的人突然就多了好多呢？

\BookQuestion{网友：腾讯put进来之后卖call，半年和一年您大概会选择什么价位？}

我一般不会想这个问题。简单讲比如一年20\%以上的回报我就满意了。价格太低我一般就懒得卖call了，但有时候也可能会卖买入价的call，因为来回的premium（期权权利金）是不错的利润哈。一般至少能满足年化15\%以上我才会动手，不然就只好闲着，哪怕只是拿着T- Bill（短期国债）也没办法。T- Bill毕竟是无风险的。（2024-10-26）
